% vim: ai sw=2 spell spelllang=en

\section{How We Do It at SUSE}
\frame{\sectionpage}

\begin{frame}{Distribution Kernel}

  \begin{itemize}[<+->]
    \item One kernel for \emph{all} users and hardware
      \begin{itemize}[<.->]
	\item Nearly everything as modules (thousands of \file{.ko})
	\item Several archs: \xesh, arm64, ppc64le, s390x, \dots
	\item Several \emph{flavors}: \file{default}, \file{rt}, \file{64kb}, \file{kvmsmall}, \dots
      \end{itemize}
    \item Reproducible: anyone can rebuild the exact same kernel
    \item Built in an isolated environment (no network, clean chroot/VM)
      \begin{itemize}[<.->]
	\item Open Build Service (OBS): \url{https://build.opensuse.org/project/subprojects/Kernel}
      \end{itemize}
    \item Tumbleweed: latest upstream stable + a few patches
    \item SLE: older base + thousands of backported patches, supported for 10+ years
  \end{itemize}

\end{frame}

\begin{frame}{\file{kernel-source.git}}

  \begin{itemize}[<+->]
    \item \emph{Not} a kernel tree! A tree of \emph{patches}
      \begin{itemize}[<.->]
	\item \url{https://github.com/SUSE/kernel-source}
	\item One branch per product: \file{stable} (Tumbleweed), \file{SL-16.1},
	  \file{SLE15-SP7}, \dots
      \end{itemize}
    \item Content
      \begin{itemize}[<.->]
	\item \file{patches.*/} + \file{series.conf} -- patches on top of upstream tarball
	\item \file{config/<arch>/<flavor>} -- \file{.config} for each build
	\item \file{supported.conf} -- support status of modules
	\item \file{rpm/*.spec.in} -- templates of RPM spec files
	\item \file{scripts/} -- tooling
      \end{itemize}
    \item \cmd{scripts/sequence-patch} -- apply patches, get a kernel tree
    \item \cmd{scripts/tar-up} -- create package sources for OBS
  \end{itemize}

\end{frame}

\begin{frame}{From \file{kernel-source.git} to \file{kernel-default.rpm}}

  \begin{center}
  \begin{tikzpicture}[node distance=6mm and 6mm]
    \node[bnode] (git) {kernel-source.git};
    \node[bnode, right=12mm of git] (pkg) {linux-x.y.tar.xz\\patches.*.tar.bz2\\
      config.tar.bz2\\kernel-default.spec\\\dots};
    \node<2->[bnode, right=12mm of pkg, fill=yellow!20] (obs) {OBS\\
      \texttt{rpmbuild}};
    \node<3->[bnode, right=10mm of obs, fill=green!20] (rpm) {kernel-default\\
      kernel-default-devel\\kernel-default-debuginfo\\kernel-devel\\kernel-syms\\\dots};
    \node<4->[bnode, below=10mm of rpm] (qa) {openQA};
    \node<4->[bnode, left=14mm of qa] (tw) {Tumbleweed\\update};

    \draw[barrow] (git) -- node[btool, above] {tar-up} (pkg);
    \draw<2->[barrow] (pkg) -- node[btool, above] {osc / push} (obs);
    \draw<3->[barrow] (obs) -- (rpm);
    \draw<4->[barrow] (rpm) -- (qa);
    \draw<4->[barrow] (qa) -- (tw);
  \end{tikzpicture}
  \end{center}

  \begin{itemize}[<+(1)->]
    \item Spec \texttt{\%prep}: unpack tarball, apply patches, copy config
    \item Spec \texttt{\%build}: \cmd{make olddefconfig}, then \cmd{make} --
      the same as today!
    \item Spec \texttt{\%install}: \cmd{make install}, \cmd{modules_install},
      split debuginfo, sign modules, \dots
  \end{itemize}

\end{frame}

\begin{frame}[fragile]{What Is in \file{kernel-default.rpm}?}

  \begin{lstlisting}[language=bash]
$ rpm -ql kernel-default | grep -v '/kernel/'
/usr/lib/modules/7.2.8-1-default/System.map
/usr/lib/modules/7.2.8-1-default/config
/usr/lib/modules/7.2.8-1-default/modules.order
/usr/lib/modules/7.2.8-1-default/vmlinuz
...
$ rpm -ql kernel-default | grep -c '\.ko'
5413
  \end{lstlisting}

  \pause

  \begin{itemize}[<+->]
    \item Scriptlets on install: \cmd{depmod}, \cmd{dracut} (initrd), bootloader entry
    \item \file{kernel-default-devel} -- for out-of-tree modules
      (\file{/lib/modules/*/build})
    \item \file{kernel-default-debuginfo} -- \file{vmlinux} with DWARF (for \cmd{crash}, \ldots)
    \item Upstream has packaging targets too
      \begin{itemize}[<.->]
	\item \cmd{make binrpm-pkg}, \cmd{rpm-pkg}, \cmd{bindeb-pkg}, \cmd{pacman-pkg}, \cmd{tar-pkg}
	\item Do not work well on SUSE
      \end{itemize}
  \end{itemize}

\end{frame}

